% concatenated from article.tex
\documentclass[a4paper,10pt]{article}

\usepackage{authblk}

\usepackage[utf8]{inputenc} % allow utf-8 input
% \usepackage[frenchb]{babel}
% \usepackage[latin1]{inputenc}
\usepackage[T1]{fontenc}    % use 8-bit T1 fonts
\usepackage{indentfirst}
\usepackage{lmodern}

\usepackage{mathtools}
\usepackage{verbatim}
\usepackage{url}            % simple URL typesetting
\usepackage{booktabs}       % professional-quality tables
\usepackage{amsfonts}       % blackboard math symbols
\usepackage{amsmath}
\usepackage{amssymb}
\usepackage{amsthm}
\usepackage{microtype}      % microtypography
\usepackage[a4paper]{geometry} % reduce margins
\usepackage{enumitem}
\usepackage{xcolor}
\usepackage{mathrsfs}

\usepackage{geometry}
%\geometry{margin=2.5cm}
\usepackage{fancyhdr}
\usepackage{listings}

\usepackage[linktocpage=true]{hyperref}
\hypersetup{
    colorlinks=true,
    linkcolor=blue,
    citecolor=violet,
    urlcolor=blue
}

\usepackage{graphicx}
\usepackage[makeroom]{cancel}
\usepackage{tikz}

\usepackage[inkscapeformat=png]{svg}
\usepackage[normalem]{ulem}


\sloppy

\theoremstyle{plain}
\newtheorem{theorem}{Theorem}[section]
\newtheorem{proposition}[theorem]{Proposition}
\newtheorem{corollary}[theorem]{Corollary}
\newtheorem{lemma}[theorem]{Lemma}

\theoremstyle{definition}
\newtheorem{definition}[theorem]{Definition}
\newtheorem{example}[theorem]{Example}  %% And a not so common one.
\newtheorem{exercise}[theorem]{Exercise}
\newtheorem{problem}[theorem]{Problem}

\theoremstyle{remark}
\newtheorem{remark}[theorem]{Remark}

% \newenvironment{proof}{{\sc Proof:}}{~\hfill QED}

\DeclareMathOperator*{\essinf}{ess\,inf}
\DeclareMathOperator*{\esssup}{ess\,sup}
\DeclareMathOperator{\Tr}{Tr}
\DeclareMathOperator{\Gram}{Gram}
\DeclareMathOperator{\Span}{span}
\DeclareMathOperator{\Dim}{dim}
\DeclareMathOperator{\Terms}{Terms}
\DeclareMathOperator{\cl}{cl}
\DeclareMathOperator{\Ker}{Ker}
\DeclareMathOperator{\MultiLinForm}{MultiLinForm}
\DeclareMathOperator{\Gr}{Gr}      


\makeatletter
\def\keywords{\xdef\@thefnmark{}\@footnotetext}
\makeatother

\newcommand*\circled[1]{\tikz[baseline=(char.base)]{\node[shape=circle,draw,inner sep=2pt] (char) {#1};}}

%%%%%%%%%%%%%%%%%%%%%%%%%%      
%%%%%%%%%%%%%%%%%%%%%%%%%%
%%%%%%%%%%%%%%%%%%%%%%%%%%
            

\title{Two Hornich-Hlawka-type and Gram matrix-based inequalities}

\author{Nizar El Idrissi and Hicham Zoubeir}


\newcommand{\Addresses}{{% additional braces for segregating \footnotesize
  \bigskip
  \footnotesize

  \textbf{Nizar El Idrissi.}
  \par\nopagebreak Laboratoire : Equations aux dérivées partielles, Algèbre et Géométrie spectrales.
  \par\nopagebreak
  Département de mathématiques, faculté des sciences, université Ibn Tofail, 14000 Kénitra.\par\nopagebreak 
  \textit{E-mail address} : \texttt{nizar.elidrissi@uit.ac.ma}
  
  \medskip
  
  \textbf{Hicham Zoubeir.}
  \par\nopagebreak Laboratoire : Equations aux dérivées partielles, Algèbre et Géométrie spectrales.
  \par\nopagebreak
  Département de mathématiques, faculté des sciences, université Ibn Tofail, 14000 Kénitra.\par\nopagebreak 
  \textit{E-mail address} : \texttt{hzoubeir2014@gmail.com}

}}

%%%%%%%%%%%%%%%%%%%%%
%%%%%%%%%%%%%%%%%%%%%
%%%%%%%%%%%%%%%%%%%%%

\begin{document}
\newpage
\maketitle

\begin{abstract}
We establish two inequalities in real inner product spaces. The first is a multiplicative strengthening of the classical Hornich–Hlawka inequality: for all vectors $x,y,z$ in a real inner product space $H$
\[ \|x\|\,\|y\| + \|z\|\,\|x+y+z\| \;\geq\; \|x+z\|\,\|y+z\|. \]
We provide a complete characterization of the equality cases in terms of the linear dependence of $x,y,z$, and explicit conditions on their Gram matrix, showing in particular that equality occurs only in flat (at most two-dimensional) configurations. We also show that this inequality implies the classical Hornich–Hlawka inequality, thereby establishing a strict hierarchy between the two. The second result is a parametric inequality derived from the positive semidefiniteness of Gram matrices: for all $x,y,z \in H$ and $\alpha, \beta, \gamma \in \mathbb{R}$, 
\[ \alpha^2\|x\|^2\langle y,z\rangle^2 + \beta^2\|y\|^2\langle x,z\rangle^2 + \gamma^2\|z\|^2\langle x,y\rangle^2 + 2(\alpha\beta + \alpha\gamma + \beta\gamma)\langle x,y\rangle\langle x,z\rangle\langle y,z\rangle \;\geq\; 0. \]
Optimizing over the parameters yields sharp inequalities relating the pairwise inner products and norms of three vectors, which can be viewed as reverse inequalities to the Gram determinant inequality $\det G \geq 0$. As a special case, this recovers and strengthens the classical Cauchy–Schwarz inequality.
\end{abstract}

\keywords{2020 \emph{Mathematics Subject Classification.} 46C05; 26D15 }
\keywords{\emph{Key words and phrases.} Hornich-Hlawka inequality; Gram matrix; Cauchy-Schwarz inequality; inner product spaces; norm inequalities }
%\def\thefootnote{*}\footnotetext{These authors contributed equally to this work.}


\tableofcontents

\section{Introduction}

The Hornich–Hlawka inequality — sometimes called Hlawka's inequality or the quadrilateral inequality — is a classical result in the geometry of normed spaces. In its standard form, it asserts that in any normed vector space $E$ in which the inequality holds,
\[\|x\| + \|y\| + \|z\| + \|x+y+z\| \;\geq\; \|x+y\| + \|x+z\| + \|y+z\| \qquad \forall\, x,y,z \in E.\]
The inequality was originally established by Hornich \cite{Hornich1942} and independently by Hlawka \cite{Hlawka1990}, and has since generated a substantial literature, being studied extensively from the perspectives of functional analysis, convex geometry, and the theory of inequalities. Fechner \cite{Fechner2014} initiated the systematic study of the associated functional inequality $f(x+y)+f(y+z)+f(x+z) \leq f(x+y+z)+f(x)+f(y)+f(z)$, investigating it for unknown mappings $f$ defined on Abelian groups and linear spaces. In \cite{Fechner2017}, he further analyzed the regularity of solutions of this functional inequality on topological abelian groups, establishing conditions under which measurability or local boundedness forces continuity, and characterizing the lower and upper hull of $f$. Fechner and Jabłońska \cite{FechnerJablonska2025} extended this line of inquiry to the almost-everywhere setting, studying the functional inequality for mappings $f:\mathbb{R}\to\mathbb{R}$ that need only be satisfied outside a Lebesgue null set. In a different direction, Honda, Okazaki and Takahashi \cite{HondaOkazakiTakahashi1998} generalized Hlawka's three-vector inequality in Hilbert space to an $n$-element inequality in $L^1$, relating it to Hanner's inequality and to the Adamović inequality. Ressel \cite{Ressel2015} showed that the Hornich–Hlawka inequality, initially stated for the identity function on the real line, extends to the entire class of Bernstein functions on the half-line, and proved in particular that it holds for the square-root function when applied to vectors in Euclidean space. Takahashi, Takahasi and Wada \cite{TakahashiTakahasiWada2009} investigated a general integral form of Hlawka's inequality on Banach spaces, establishing both the inequality and its reverse and clarifying the connection with the Djoković inequality in a weighted integral setting. Serre \cite{Serre2015} considered a genuinely different geometric setting: he proved that in the future cone of Minkowski space, the Lorentzian pseudo-norm satisfies a "reverse" Hlawka-type inequality — the orientation of the inequality is opposite to the Euclidean case, in direct analogy with the reverse of the Cauchy–Schwarz inequality that holds in that setting. While the inequality fails in certain normed spaces, it is known to hold in all inner product spaces, as well as in broader classes including $L^p$ spaces for $p \in [1,2]$. \\ \\
The other inequality discussed in the present paper belongs to the tradition of inequalities in inner product spaces derived from structural properties of Gram matrices. The nonnegativity of the Gram determinant $\det G \geq 0$ and the positive semidefiniteness of Gram matrices are classical tools for generating vector inequalities; see, e.g., \cite{AravHallLi2008}, where a Cauchy-Schwarz-type inequality for triples of vectors is established and interpreted geometrically in terms of the vertex angles of a tetrahedron, with Gram matrices playing a central role. A major impetus for the systematic study of such inequalities is provided by Buzano's inequality and its relatives. Dragomir \cite{Dragomir2005} established refinements of both Buzano's inequality and Kurepa's inequality in inner product spaces, using Gramian-based arguments to obtain improvements of Grüss-type inequalities and reverse Schwarz inequalities. Sababheh, Moradi and Heydarbeygi \cite{SababhehMoradiHeydarbeygi2022} refined the Cauchy-Schwarz and Buzano inequalities simultaneously and provided a new proof of a refined Kreĭn-type inequality, with applications to operator norm and numerical radius inequalities for Hilbert space operators. For a comprehensive treatment of the broader landscape of Schwarz-type inequalities in inner product spaces, including reverses, refinements, and connections to Grüss and Bessel inequalities, we refer to the monograph \cite{Dragomir2005book}. The inequality established in Section~\ref{section-3} of the present paper can be situated within this framework: it derives from the positive semidefiniteness of the Gram matrix via a single substitution, and its corollaries constitute sharp reverse inequalities to the Gram determinant condition, recovering the Cauchy-Schwarz inequality as a special case. \\ \\
The present paper makes two independent contributions to the theory of inequalities in inner product spaces. \\ \\
\textbf{First contribution.} In Section~\ref{section-2}, we prove that in any real inner product space $H$, the Hornich–Hlawka inequality admits the following multiplicative strengthening: for all $x,y,z \in H$,
\[\|x\|\,\|y\| + \|z\|\,\|x+y+z\| \;\geq\; \|x+z\|\,\|y+z\|.\]
This inequality is genuinely stronger than the classical one: we show that it implies the Hornich–Hlawka inequality via a direct algebraic argument, establishing a strict hierarchy. The proof proceeds by reducing the inequality to the nonnegativity of a quartic form $\xi$ in the entries of the Gram matrix of $x,y,z$. Viewing $\xi$ as a function over the convex cone of $3\times 3$ positive semidefinite matrices, we show that its minimum over this cone is attained on the boundary — that is, in configurations where the Gram matrix is singular and the three vectors are linearly dependent. The verification that $\xi \geq 0$ on this boundary is then carried out by direct substitution and algebraic factorization. We also provide a complete description of the equality cases (Proposition~\ref{proposition-equality-case}), showing that equality holds if and only if $x,y,z$ are linearly dependent and satisfy explicit algebraic conditions on their norms and inner products. In particular, equality can only occur in flat configurations, i.e., when the three vectors lie in a subspace of dimension at most two. \\ \\
\textbf{Second contribution.} In Section~\ref{section-3}, we prove a parametric family of inequalities that follows directly from the positive semidefiniteness of the Gram matrix $G$ associated to three vectors $x,y,z \in H$. Namely, we show that for all $\alpha, \beta, \gamma \in \mathbb{R}$,
\[\alpha^2\|x\|^2\langle y,z\rangle^2 + \beta^2\|y\|^2\langle x,z\rangle^2 + \gamma^2\|z\|^2\langle x,y\rangle^2 + 2(\alpha\beta + \alpha\gamma + \beta\gamma)\langle x,y\rangle\langle x,z\rangle\langle y,z\rangle \;\geq\; 0.\]
The proof is elementary: it consists in substituting $(\alpha\langle y,z\rangle,\, \beta\langle x,z\rangle,\, \gamma\langle x,y\rangle)$ as a test vector in the defining quadratic form of $G$. The content lies rather in the corollaries obtained by optimizing the parameters $\alpha, \beta, \gamma$. When $\langle x,y\rangle\langle x,z\rangle\langle y,z\rangle > 0$, the choice $(\alpha,\beta,\gamma) = (1,1,-2)$ (and its cyclic permutations) yields the sharp inequality 
\[ \|x\|^2\langle y,z\rangle^2 + \|y\|^2\langle x,z\rangle^2 + \|z\|^2\langle x,y\rangle^2 \;\geq\; 3\langle x,y\rangle\langle x,z\rangle\langle y,z\rangle, \] 
while when $\langle x,y\rangle\langle x,z\rangle\langle y,z\rangle < 0$, the choice $(\alpha,\beta,\gamma) = (1,1,1)$ yields the sharp inequality
\[\|x\|^2\langle y,z\rangle^2 + \|y\|^2\langle x,z\rangle^2 + \|z\|^2\langle x,y\rangle^2 \;\geq\; -6\langle x,y\rangle\langle x,z\rangle\langle y,z\rangle.\]
Both inequalities are sharp, as demonstrated by explicit examples. These results can be understood as reverse inequalities to the Gram determinant condition $\det G \geq 0$, and the first of them implies the Cauchy–Schwarz inequality as the special case $z = y$. \\
\mbox{} \\
\textbf{Notation.} Throughout, $H$denotes a real inner product space, with inner product $\langle \cdot, \cdot \rangle$ and associated norm $\|\cdot\| = \sqrt{\langle \cdot, \cdot \rangle}$.

\section{A stronger Hornich-Hlawka inequality}
\label{section-2}

\begin{theorem}[Strong Hornich–Hlawka inequality]
For all $x,y,z \in H$, we have
\[ \lVert x \rVert \lVert y \rVert + \lVert z \rVert \lVert x+y+z \rVert \geq \lVert x + z \rVert \lVert y + z \rVert, \]
Moreover, equality holds if and only if $x,y,z$ are linearly dependent and their dependence coefficients satisfy one of the equivalent conditions listed in proposition \ref{proposition-equality-case}
\end{theorem}

\begin{remark} \mbox{}\\*
\begin{itemize}
\item When $H = \mathbb{R}$ or $\mathbb{C}$, the inequality can be easily proved by expanding both sides and using the multiplicativity of the modulus (this property is senseless when we consider vector norms instead of the modulus of complex numbers).
\item This stronger Hornich-Hlawka inequality implies the classical Hornich-Hlawka inequality. Indeed, since both sides of the classical Hornich-Hlawka inequality are nonnegative, we may replace them with their squares. \\
For the left-hand side, we obtain
\[
\bigl( \|x\| + \|y\| + \|z\| + \|x + y + z\| \bigr)^2 = 2\bigl(\|x\|^2 + \|y\|^2 + \|z\|^2\bigr) + 2\bigl(\langle x,y\rangle + \langle y,z\rangle + \langle z,x\rangle\bigr)
\]
\[
+ 2\bigl( \|x\| + \|y\| + \|z\| \bigr) \|x + y + z\| + 2\bigl( \lVert x \rVert \lVert y \rVert + \lVert x \rVert \lVert z \rVert + \lVert y \rVert \lVert z \rVert \bigr).
\]
For the right-hand side, we have
\[
\bigl( \|x + y\| + \|y + z\| + \|z + x\| \bigr)^2 = 2\bigl(\|x\|^2 + \|y\|^2 + \|z\|^2\bigr) + 2\bigl(\langle x,y\rangle + \langle y,z\rangle + \langle z,x\rangle\bigr)
\]
\[
+ 2\bigl( \|x + y\| \cdot \|y + z\| + \|y + z\| \cdot \|z + x\| + \|z + x\| \cdot \|x + y\| \bigr).
\]

Thus, the classical Hornich-Hlawka inequalityis is equivalent to
\[
\bigl( \|x\| + \|y\| + \|z\| \bigr) \|x + y + z\| + \bigl( \lVert x \rVert \lVert y \rVert +  \lVert x \rVert \lVert z  \rVert + \lVert y \rVert \lVert z \rVert \bigr).
\]
\[
\geq \|x + y\| \cdot \|y + z\| + \|y + z\| \cdot \|z + x\| + \|z + x\| \cdot \|x + y\|
\]
which, after regrouping, is equivalent to
\[
\bigl( \lVert x \rVert \lVert y \rVert  + \|z\| \cdot \|x + y + z\| - \|y + z\| \cdot \|z + x\| \bigr) +
\]
\[
\bigl( \lVert y \rVert \lVert z \rVert + \|x\| \cdot \|x + y + z\| - \|z + x\| \cdot \|x + y\| \bigr) +
\]
\[
\bigl( \lVert x \rVert \lVert z  + \|y\| \cdot \|x + y + z\| - \|x + y\| \cdot \|y + z\| \bigr) \geq 0.
\]
Therefore, it is clear that the stronger Hornich-Hlawka inequality implies the classical Hornich-Hlawka inequality.
\end{itemize}
\end{remark}

\begin{proof} \mbox{}\\* \mbox{}\\*
\textbf{Idea:} we reduce the inequality to a quartic form in the Gram matrix entries and show that its minimum over the PSD cone is attained on the boundary. \\
\begin{itemize}
\item \textbf{First step. Reduction.} \\
Denote by $L$ and $R$ the left and right hand sides of this inequality, respectively. \\
We have
\begin{align*}
L^2 &= \lVert x \rVert ^2 \lVert y \rVert^2 + \lVert z \rVert^2 \left( \lVert x \rVert^2 + \lVert y \rVert^2 + \lVert z \rVert^2 + 2 \langle x,y \rangle + 2 \langle x,z \rangle + 2\langle y,z \rangle  \right) + 2 \lVert x \rVert \lVert y \rVert \lVert z \rVert \lVert x + y + z \rVert \\
&=  \lVert x \rVert ^2 \lVert y \rVert^2 + \lVert z \rVert^2 \lVert x \rVert^2 + \lVert z \rVert^2 \lVert y \rVert^2 + \lVert z \rVert^4 + 2\lVert z \rVert^2 \langle x , y \rangle +  2\lVert z \rVert^2 \langle x , z \rangle +  2\lVert z \rVert^2 \langle y , z \rangle \\
&\quad \quad + 2 \lVert x \rVert \lVert y \rVert \lVert z \rVert \lVert x + y + z \rVert
\end{align*}
and
\begin{align*}
R^2 &= \left( \lVert x \rVert^2 + \lVert z \rVert^2 + 2\langle x , z \rangle \right) \left( \lVert y \rVert^2 + \lVert z \rVert^2 + 2\langle y , z \rangle \right) \\
&= \lVert x \rVert ^2 \lVert y \rVert^2 + \lVert z \rVert^2 \lVert x \rVert^2 + \lVert z \rVert^2 \lVert y \rVert^2 + \lVert z \rVert^4  + 2 \langle x , z \rangle \lVert y \rVert^2 + 2\langle x , z \rangle \lVert z \rVert^2 \\
&\quad \quad + 2 \langle y , z \rangle \lVert x \rVert^2 + 2 \langle y , z \rangle \lVert z \rVert^2 + 4 \langle x , z \rangle \langle y , z \rangle
\end{align*}
We have
\begin{align*}
&L \geq R \\
&\Leftrightarrow L^2 \geq R^2 \\
&\Leftrightarrow \frac{L^2 - \left( \lVert x \rVert ^2 \lVert y \rVert^2 + \lVert z \rVert^2 \lVert x \rVert^2 + \lVert z \rVert^2 \lVert y \rVert^2 + \lVert z \rVert^4 \right)}{2} \geq \frac{R^2 - \left( \lVert x \rVert ^2 \lVert y \rVert^2 + \lVert z \rVert^2 \lVert x \rVert^2 + \lVert z \rVert^2 \lVert y \rVert^2 + \lVert z \rVert^4 \right)}{2} \\
&\Leftrightarrow \lVert z \rVert^2 \langle x , y \rangle +  \lVert z \rVert^2 \langle x , z \rangle +  \lVert z \rVert^2 \langle y , z \rangle + \lVert x \rVert \lVert y \rVert \lVert z \rVert \lVert x + y + z \rVert \geq \\
&\quad \quad \quad \langle x , z \rangle \lVert y \rVert^2 + \langle x , z \rangle \lVert z \rVert^2 +  \langle y , z \rangle \lVert x \rVert^2 + \langle y , z \rangle \lVert z \rVert^2 + 2 \langle x , z \rangle \langle y , z \rangle \\
&\Leftrightarrow \lVert z \rVert^2 \langle x , y \rangle + \lVert x \rVert \lVert y \rVert \lVert z \rVert \lVert x + y + z \rVert \geq \lVert x \rVert^2 \langle y , z \rangle + \lVert y \rVert^2 \langle x , z \rangle + 2 \langle x , z \rangle \langle y , z \rangle
\end{align*}
This inequality is equivalent to
\[ \Leftrightarrow \lVert x \rVert \lVert y \rVert \lVert z \rVert \lVert x + y + z \rVert \geq \lVert x \rVert^2 \langle y , z \rangle + \lVert y \rVert^2 \langle x , z \rangle - \lVert z \rVert^2 \langle x , y \rangle + 2 \langle x , z \rangle \langle y , z \rangle \]
Denote by $\textbf{L}$ and $\textbf{R}$ the left and right hand sides of this last inequality, respectively. \\
If $\textbf{R} \leq 0$, the inequality is trivial, so we can suppose that $\textbf{R} \geq 0$. \\
Set 
\[ a = \lVert x \rVert, \quad b = \lVert y \rVert, \quad c = \lVert z \rVert, \quad p = \langle x,y \rangle, \quad q = \langle x,z \rangle, \quad r = \langle y,z \rangle. \]
We have 
\begin{align*}
\textbf{L}^2 &=  \lVert x \rVert^2  \lVert y \rVert^2  \lVert z \rVert^2  \left( \lVert x \rVert^2 + \lVert y \rVert^2  + \lVert z \rVert^2 + 2  \langle x , y \rangle + 2 \langle x , z \rangle +  2 \langle y , z \rangle \right) \\
&= a^2 b^2 c^2 (a^2 + b^2 + c^2 + 2p + 2q + 2r).
\end{align*}
and
\begin{align*}
\textbf{R}^2 &= \lVert x \rVert^4 \langle y , z \rangle^2 +  \lVert y \rVert^4 \langle x , z \rangle^2 +  \lVert z \rVert^4 \langle x , y \rangle^2 +  4 \langle x , z \rangle^2 \langle y , z \rangle^2 \\
&\quad + 2 \lVert x \rVert^2 \lVert y \rVert^2 \langle x , z \rangle \langle y , z \rangle - 2 \lVert x \rVert^2 \lVert z \rVert^2  \langle x , y \rangle  \langle y , z \rangle   - 2 \lVert y \rVert^2 \lVert z \rVert^2 \langle x , y \rangle  \langle x , z \rangle   \\
&\quad + 4 \lVert x \rVert^2  \langle x , z \rangle \langle y , z \rangle ^2 + 4 \lVert y \rVert^2 \langle y , z \rangle \langle x , z \rangle ^2 - 4 \lVert z \rVert^2 \langle x , y \rangle \langle x , z \rangle  \langle y , z \rangle  \\
&=  a^4 r^2 + b^4 q^2 + c^4 p^2 + 4q^2 r^2 + 2a^2 b^2 qr - 2a^2 c^2 pr - 2b^2 c^2 pq + 4a^2 q r^2 + 4b^2 r q^2 - 4c^2 pqr
\end{align*}
Consider 
\begin{align*}
\xi &= \textbf{L}^2 - \textbf{R}^2 \\
&= a^2 b^2 c^2 (a^2 + b^2 + c^2 + 2p + 2q + 2r) \\
&\quad - \left( a^4 r^2 + b^4 q^2 + c^4 p^2 + 4q^2 r^2 + 2a^2 b^2 qr - 2a^2 c^2 pr - 2b^2 c^2 pq + 4a^2 q r^2 + 4b^2 r q^2 - 4c^2 pqr \right)
\end{align*}

\item \textbf{Second step. The minimum of $\xi$ over the PSD cone is attained on the boundary.} \\
Consider $\xi$ as a quadratic function in one variable, say $p$, for fixed $a, b, c, q, r$ (satisfying the relevant conditions, see next). \\ 
The coefficient of $p^2$ is $-c^4 \leq 0$, so it is a downward-opening parabola. \\
The set of allowed values $a, b, c \geq 0$, $p, q, r$ are such that the Gram matrix
\[
G = \begin{pmatrix}
a^2 & p & q \\
p & b^2 & r \\
q & r & c^2
\end{pmatrix}
\]
is positive semidefinite. \\
This requires all principal minors to be nonnegative: $|p| \leq ab$, $|q| \leq ac$, $|r| \leq bc$, and $\det G \geq 0$. \\
Since the allowed range for $p$ is the interval where $\det G \geq 0$ and $|p| \leq ab$, which is closed and bounded, the minimum of $\xi$ over this interval occurs at one of the endpoints. \\
The endpoints correspond to boundaries where either $\det G = 0$ or $p = \pm ab$ (but note that $p = \pm ab$ implies the 2x2 minor for $x,y$ is zero, which implies $x,y$ are linearly dependent, which forces $\det G = 0$). Thus, the minima occur on the boundary where $G$ is singular ($\det G = 0$). \\
To show $\xi \geq 0$ on this boundary, assume linear dependence:
\begin{itemize}
\item If $z = \lambda x + \mu y$, then substitute
\begin{align*}
c^2 &= \lambda^2 a^2 + \mu^2 b^2 + 2 \lambda \mu p \\
q &= \lambda a^2 + \mu p \\
r &= \lambda p + \mu b^2
\end{align*}
After substitution and simplification, $\xi$ factors as
\[ \xi = (a^2 b^2 - p^2) \left( a^2 \lambda (\lambda + 1) - b^2 \mu (\mu + 1) \right)^2. \]
The first factor is nonnegative by the Cauchy-Schwarz inequality, and the second factor is a square, hence nonnegative. Therefore, $\xi \geq 0$.

\item If $x = \lambda y + \mu z$, then substitute
\begin{align*}
a^2 &= \lambda^2 b^2 + \mu^2 c^2 + 2 \lambda \mu r, \\
p &= \lambda b^2 + \mu r, \\
q &= \lambda r + \mu c^2.
\end{align*}
After substitution and simplification, $\xi$ factors as
\begin{equation*}
\xi  = \Bigl( b^2 c^2 - r^2 \Bigr)
       \Bigl( b^2 \lambda (\lambda + 1) + c^2 \mu (\mu + 1) + 2 \lambda (\mu + 1) r \Bigr)^2.
\end{equation*}
The first factor is nonnegative by the Cauchy-Schwarz inequality, and the second factor is a square, hence nonnegative. Therefore, $\xi \geq 0$.
\item If $y = \lambda x + \mu z$, then substitute
\begin{align*}
b^2 &= \lambda^2 a^2 + \mu^2 c^2 + 2 \lambda \mu q, \\
p &= \lambda a^2 + \mu q, \\
r &= \lambda q + \mu c^2.
\end{align*}
After substitution and simplification, $\xi$ factors as
\[
\xi = \left( a^2 c^2 - q^2 \right) \left( a^2 \lambda (\lambda + 1) + c^2 \mu (\mu + 1) + 2 \lambda (\mu + 1) q \right)^2.
\]
The first factor is nonnegative by the Cauchy-Schwarz inequality, and the second factor is a square, hence nonnegative. Therefore, $\xi \geq 0$.
\end{itemize}
Therefore, $\xi \geq 0$ on the boundary. 

\item \textbf{Conclusion.} \\
Since $\xi \geq 0$ on the boundary, we have $\xi \geq 0$ throughout the domain. \\
Equality arises in configurations where the vectors are linearly dependent, with additional conditions on the scalar products and norms: the inequality becomes an equality when the geometry is "flat" (vectors in a plane or a line).
\end{itemize}
\end{proof}

\begin{proposition}[Equality cases]
\label{proposition-equality-case}
Let $x,y,z\in H$.
Equality holds in the inequality
\[
\|x\|\|y\|+\|z\|\|x+y+z\|=\|x+z\|\|y+z\|
\]
if and only if the vectors $x,y,z$ are linearly dependent and one of the following cases occurs.

\begin{enumerate}
\item[\textup{(i)}]
There exist real numbers $\lambda,\mu$ such that
\[
z=\lambda x+\mu y
\]
and
\[
\|x\|^{2}\,\lambda(\lambda+1)
=
\|y\|^{2}\,\mu(\mu+1).
\]

\item[\textup{(ii)}]
There exist real numbers $\lambda,\mu$ such that
\[
x=\lambda y+\mu z
\]
and
\[
\|y\|^{2}\,\lambda(\lambda+1)
+
\|z\|^{2}\,\mu(\mu+1)
+
2\lambda(\mu+1)\langle y,z\rangle
=0.
\]

\item[\textup{(iii)}]
There exist real numbers $\lambda,\mu$ such that
\[
y=\lambda x+\mu z
\]
and
\[
\|x\|^{2}\,\lambda(\lambda+1)
+
\|z\|^{2}\,\mu(\mu+1)
+
2\lambda(\mu+1)\langle x,z\rangle
=0.
\]
\end{enumerate}

In all cases, the Gram matrix
\[
G=
\begin{pmatrix}
\|x\|^2 & \langle x,y\rangle & \langle x,z\rangle\\
\langle x,y\rangle & \|y\|^2 & \langle y,z\rangle\\
\langle x,z\rangle & \langle y,z\rangle & \|z\|^2
\end{pmatrix}
\]
is singular. In particular, equality occurs only in flat configurations, i.e.,
when the vectors $x,y,z$ lie in a common one or two dimensional subspace.
\end{proposition}


\section{An inequality based on Gram matrices}
\label{section-3}

\begin{theorem}
We have, for all vectors $x,y,z \in H$ and real numbers $\alpha,\beta,\gamma$:
\[ \alpha^2  \|x\|^2 \langle y, z \rangle^2 + \beta^2 \|y\|^2 \langle x, z \rangle^2 + \gamma^2 \|z\|^2 \langle x, y \rangle^2  + 2  (\alpha\beta + \alpha\gamma + \beta\gamma) \langle x, y\rangle \langle x, z\rangle \langle y, z\rangle \geq 0. \]
\end{theorem}

\begin{proof}
Given \(x, y, z\) in $H$, the Gram matrix associated to $(x,y,z)$ is  
\[
G = \begin{pmatrix}
\langle x, x \rangle & \langle x, y \rangle & \langle x, z \rangle \\
\langle y, x \rangle & \langle y, y \rangle & \langle y, z \rangle \\
\langle z, x \rangle & \langle z, y \rangle & \langle z, z \rangle
\end{pmatrix}.
\]
It is positive semidefinite, so for any real numbers \(\alpha, \beta, \gamma\), we have  
\[
Q(\alpha, \beta, \gamma) = (\alpha, \beta, \gamma) \, G \, \begin{pmatrix} \alpha \\ \beta \\ \gamma \end{pmatrix} \ge 0.
\]
Expanding, we have
\[
Q(\alpha, \beta, \gamma) =
\alpha^2 \|x\|^2 + \beta^2 \|y\|^2 + \gamma^2 \|z\|^2
+ 2\alpha\beta\langle x, y\rangle + 2\alpha\gamma\langle x, z\rangle + 2\beta\gamma\langle y, z\rangle \geq 0.
\]
Replace
\[
\alpha \mapsto \alpha \langle y, z \rangle,
\]
\[
\beta \mapsto \beta \langle x, z \rangle,
\]
\[
\gamma \mapsto \gamma \langle x, y \rangle.
\]
Then
\begin{align*}
R(\alpha,\beta,\gamma) &= Q(\alpha \langle y, z \rangle, \beta \langle x, z \rangle , \gamma  \langle x, y \rangle) \\
&=  \alpha^2  \|x\|^2 \langle y, z \rangle^2 + \beta^2 \|y\|^2 \langle x, z \rangle^2 + \gamma^2 \|z\|^2 \langle x, y \rangle^2  \\
&\quad + 2  (\alpha\beta + \alpha\gamma + \beta\gamma) \langle x, y\rangle \langle x, z\rangle \langle y, z\rangle \geq 0.
\end{align*}
\end{proof}

\begin{corollary} \mbox{}\\*
\begin{itemize}
\item If $\langle x, y\rangle \langle x, z\rangle \langle y, z\rangle > 0$, we have
\[ \|x\|^2 \langle y, z \rangle^2 + \|y\|^2 \langle x, z \rangle^2 + \|z\|^2 \langle x, y \rangle^2  \geq 3 \langle x, y\rangle \langle x, z\rangle \langle y, z\rangle. \]
\item If $\langle x, y\rangle \langle x, z\rangle \langle y, z\rangle < 0$, we have
\[ \|x\|^2 \langle y, z \rangle^2 + \|y\|^2 \langle x, z \rangle^2 + \|z\|^2 \langle x, y \rangle^2  \geq - 6 \langle x, y\rangle \langle x, z\rangle \langle y, z\rangle.\]
\end{itemize}
These inequalities are \textbf{sharp}.
\end{corollary}

\begin{remark} \mbox{}\\* 
\begin{itemize}
\item This corollary can be understood as a reverse inequality of the Gram determinant inequality
\[ \det G \geq 0 \]
which translates to
\[ \|x\|^2 \|y\|^2 \|z\|^2 + 2\langle x, y \rangle \langle x, z \rangle \langle y, z \rangle \geq \|x\|^2 \langle y, z \rangle^2 + \|y\|^2 \langle x, z \rangle^2 + \|z\|^2 \langle x, y \rangle^2. 
\]
\item The first point of the corollary implies the Cauchy-Schwarz inequality, and therefore all other inequalities deduced from it. One can show this by taking $z=y$.
\end{itemize}
\end{remark}

\begin{proof} \mbox{}\\*
\begin{itemize}
\item If 
\[ \langle x, y\rangle \langle x, z\rangle \langle y, z\rangle \geq 0 \]
Then $R(1,1,-2) > 0$ implies
\[ \|x\|^2 \langle y, z \rangle^2 + \|y\|^2 \langle x, z \rangle^2 + 4 \|z\|^2 \langle x, y \rangle^2  \geq 6 \langle x, y\rangle \langle x, z\rangle \langle y, z\rangle.\]
Summing up symmetrically and dividing by 6, we get
\[ \|x\|^2 \langle y, z \rangle^2 + \|y\|^2 \langle x, z \rangle^2 + \|z\|^2 \langle x, y \rangle^2  \geq 3 \langle x, y\rangle \langle x, z\rangle \langle y, z\rangle.\]
\item If
\[ \langle x, y\rangle \langle x, z\rangle \langle y, z\rangle < 0 \]
Then $R(1,1,1) \geq 0$ implies
\[ \|x\|^2 \langle y, z \rangle^2 + \|y\|^2 \langle x, z \rangle^2 + \|z\|^2 \langle x, y \rangle^2  \geq - 6 \langle x, y\rangle \langle x, z\rangle \langle y, z\rangle. \]
\end{itemize}
\textbf{Sharpness.} For $D>0$, take \( x = y = z \) nonzero. It yields LHS = RHS. \\
For $D < 0$, take \(x = (1,0,0)\), \(y = \left(\frac12, \frac{\sqrt{3}}{2}, 0\right)\), \(z = \left(\frac12, -\frac{\sqrt{3}}{2}, 0\right)\). \\ 
We have
\[
\langle x, y \rangle = 1\cdot \frac12 + 0\cdot \frac{\sqrt{3}}{2} = \frac12
\]
\[
\langle x, z \rangle = 1\cdot \frac12 + 0\cdot \left(-\frac{\sqrt{3}}{2}\right) = \frac12
\]
\[
\langle y, z \rangle = \frac12\cdot \frac12 + \frac{\sqrt{3}}{2} \cdot \left(-\frac{\sqrt{3}}{2}\right) = \frac14 - \frac{3}{4} = -\frac12
\]
So 
\[ RHS = -6\left(\frac12\right)\left(\frac12\right)\left(-\frac12\right) = \frac68 = \frac34. \]
Moreover
\[
\|x\|^2 \langle y, z\rangle^2 = 1\cdot \frac14 = \frac14
\]
\[
\|y\|^2 \langle x, z\rangle^2 = 1\cdot \frac14 = \frac14
\]
\[
\|z\|^2 \langle x, y\rangle^2 = 1\cdot \frac14 = \frac14
\]
So 
\[ LHS = \frac14 + \frac14 + \frac14  = \frac 34. \]
Equality holds.
\end{proof}

%%%%%%%%%%%%%%%%%%%%
\nocite{*} % include non-cited references
\bibliographystyle{amsplain}
\bibliography{references}
%%%%%%%%%%%%%%%%%%%%

\Addresses

\end{document}
